\documentclass[../tex_src/CaseBriefs.tex]{subfiles}
\usepackage{lawbrief}
\begin{document}
\NewBrief{subject={Breach of Contract, Consideration},
  plaintiff={1464-Eight, Ltd},
  defendant={Joppich},
  citation={154 S.W.3d 101 (2004)},
  course={Contracts},
  facts={The two parties entered an option contract in which Plaintiff (1464) sold Defendant (Joppich) a piece of land with the option for Plaintiff to purchase the property back at 90\% the price if Defendant failed to build a residential property on the land within 18 months.\ The contract stated that consideration for the option was a \$10 fee to be paid by Plaintiff to Defendant.\ However, that fee had never been paid (as is at least somewhat common).\ When, after 18 months, Defendant had failed to build a property, Plaintiff tried to exercise this option Defendant refused citing the failure to pay the consideration fee.},
  procedure={Defendant requests Declaratory Judgement against Plaintiff (Failed)\\
Plaintiff files counter-claim\\
Trial Court: In favor of Plaintiff\\
Appeals: In favor of Defendant\\
Supreme: In favor of Plaintiff (Reversed and remanded)},
  issue={Should the Second Restatements input on considerations be adopted in Texas?\\
Does false recital constitute a valid consideration for an option contract?},
  holding={No and Yes; Reversed and Remanded},
  principle={False recital does constitute consideration for an otherwise valid and enforcable contract.},
  reasoning={},
  opinions={Chief Justice Jefferson: Concurring; The idea that false recital is enough to constitute consideration should indicate that consideration might be a thing of the past.\ If not including a clause in the contract},
  label={case:1464-Eight,LtdVJoppich},
  notes={}
}
\end{document}